%!TeX root=MemoriaTFG.tex

\chapter{Conclusions and future work}\label{chap:conclusions}

\section{Conclusions}

Wave Launcher grew out of a simple observation: setting up a private Minecraft server still requires more technical knowledge than many players have or want to acquire. The project responds to that problem with a local Windows application that coordinates the server software, Java runtimes, Forge or Fabric, mods from Modrinth or CurseForge, configuration files, process execution and router port mappings. The result is not intended to be a general hosting panel. It is a focused desktop workflow for players who want to self-host without learning a separate administration procedure for every component.

The implementation demonstrates that the selected components can be represented behind a common application model despite substantial differences between their providers. Mojang and Adoptium distribute Java in different forms. Forge exposes Maven metadata while Fabric provides its own catalogue endpoints. Modrinth and CurseForge use different identifiers, authentication rules and response schemas. Mapping them into domain types makes server creation and migration an application operation rather than a sequence of unrelated downloads.

Hexagonal architecture was valuable during this work because external integrations could be implemented and revised independently. The later development of the interface did require service changes, but those changes occurred through explicit contracts rather than direct dependencies from pages to files or HTTP endpoints. JSON persistence also proved proportionate to a single-user desktop application, while repository interfaces retain a path toward stronger storage if future scope demands it.

The iterative weekly process suited a project whose uncertainty lay mainly in external interfaces and a new UI framework. The high-level schedule remained stable through functional completion. The final correction phase grew by one week due to MAUI form behavior, illustrating why short milestones and time reserved for integration were necessary.

The functional validation strengthens this conclusion beyond successful implementation alone. The same end-to-end scenario completed on Windows 10 Home, Windows 11 Home and Windows 11 Pro. It covered managed Java 21 and 25 installations, automatic runtime selection, Minecraft 1.21 and 26.1 servers, Forge and Fabric, CurseForge and Modrinth, dependency resolution, loader switching, a downgrade to Minecraft 26, property persistence, execution, port-conflict detection and complete cleanup. No reproducible error was observed during that procedure.

The comparison with three users also provided initial evidence for the usability objective. Manual creation took between 18:23 and 22:14, whereas the Wave Launcher workflow took between 02:08 and 05:00. Mean completion time fell from 19:45 to 03:44. During manual setup, all three users struggled with port forwarding or needed specific guidance, and interpreting successful startup and loader or mod compatibility also caused repeated uncertainty. Wave Launcher automated these responsibilities, leaving mainly local interface questions such as where to add a mod or what execution flags mean. The sample is too small to support broad statistical conclusions, but the consistent result across three different backgrounds indicates that the application reduced both completion time and the technical knowledge required in the observed task.

\section{Future work}

\subsection{Extend the empirical evaluation}

The initial comparison involved only three participants and always presented the manual task first. A larger evaluation should balance task order, use a more detailed intervention protocol and include participants from a wider range of ages and technical backgrounds. It should record completion, time, interventions, errors and user confidence consistently. This would show whether the observed reduction is maintained beyond the small exploratory sample and would separate the effect of Wave Launcher from learning during the manual task. The functional matrix should also be extended as new Minecraft, Java, Forge and Fabric versions are published, with failures classified between application defects, provider outages and unsupported router environments.

\subsection{Additional platforms}

.NET MAUI and the internal boundaries leave room for macOS or Linux desktop support, but the present product should not be described as cross-platform until adapters and complete tests exist. New process-executor and device-information implementations are required. File permissions, executable discovery, archive formats and platform packaging must also be addressed. Android and iOS are not natural server-hosting targets even though MAUI can compile interface projects for them, so desktop platforms have higher priority.

\subsection{Modpack support}

The current mod integrations manage individual projects obtained from Modrinth and CurseForge. A future version could also browse and install modpacks published by these already supported suppliers. Installing a modpack would require the application to interpret its manifest, obtain the specified loader and Minecraft version, download its mods and configuration files, and present any optional components to the user. Reusing the existing supplier adapters and compatibility model would keep modpack installation within the same server-management workflow instead of requiring users to assemble a published pack manually.

\subsection{World migration with MCA Selector}

Updating the server software does not update terrain that has already been generated in its world. Integrating MCA Selector could make world migration more useful by identifying and deleting generated chunks that have not been modified. When those areas are explored again, Minecraft would regenerate them using the terrain and features of the newer version. The launcher should present the affected area clearly and require explicit confirmation before changing the world files, since deleting the wrong chunks could remove player changes.

\subsection{Whitelist management}

Wave Launcher could provide an interface for managing the server whitelist without requiring the user to edit files or enter commands in the server console. The interface should list the currently allowed players and support adding or removing entries from the launcher. This would make access control consistent with the existing graphical management of properties, mods and execution while keeping the underlying Minecraft whitelist synchronized.

\subsection{Network diagnostics}

UPnP and NAT-PMP cover many home routers but cannot solve carrier-grade NAT, disabled discovery or firewall rules. A diagnostic screen could explain the local address, mapping protocol, public port and failure stage. Optional firewall-rule assistance and a reachability check from outside the local network would make the result easier to verify. These features must be designed cautiously because network configuration changes have security implications.

\subsection{Distribution and community development}

Wave Launcher's source code is publicly available at \url{https://github.com/nvl-ez/Wave}. The project is intended to be distributed as free software under a GNU copyleft licence.
